\documentclass[11pt,twoside]{article}

\usepackage{fancyhdr}
\usepackage{epsfig}
\usepackage{latexsym}
\usepackage{microtype}
\usepackage{lastpage} 
\usepackage{hyperref}

\usepackage[dvipsnames*,svgnames]{xcolor}
\hypersetup{colorlinks=true,allcolors=DarkBlue,linkcolor=black}

%Please include other packages if you need them.

% if you use BibTeX
%\bibliographystyle{plainurl}
% please also provide DOI information in your BibTeX entries

%------------------------------------------------------------------------------
% Please do not modify these values
\setlength{\textwidth}{12.6cm}
\setlength{\textheight}{19cm}
\setlength{\evensidemargin}{1.9cm}
\setlength{\oddsidemargin}{1.9cm}
%------------------------------------------------------------------------------

\newtheorem{definition}{Definition}
\newtheorem{theorem}{Theorem}
\newtheorem{proposition}{Proposition}
\newtheorem{corollary}{Corollary}
\newtheorem{lemma}{Lemma}
\newtheorem{remark}{Remark}
%Please define other environments if you need them

\newcommand{\proof}[1]{\par\noindent {\bf Proof:}\ \ \ #1 \hfill$\Box$\hspace{1ex}}
\newcommand{\sketch}{\par\noindent {\bf Sketch of the proof:} \hspace{1ex}}

%\setcounter{page}{}

%------------------------------------------------------------------------------
% For headings

\thispagestyle{empty}

\newcommand{\TheAuthor}{}
\newcommand{\Author}[1]{\renewcommand{\TheAuthor}{#1}}

\newcommand{\TheTitle}{}
\newcommand{\Title}[1]{\renewcommand{\TheTitle}{#1}}

\Author{Raziel Carvajal-G{\'o}mez}
\Title{Towards an Adaptive Approach for Broadcast Algorithms in MANETs}

\fancyfoot{} % clear all footer fields

%------------------------------------------------------------------------------

\begin{document}

\pagenumbering{gobble}
\vspace{0.2cm}

\begin{center}
{\Large\bf Towards an Adaptive Approach for Broadcast Algorithms in MANETs}
\end{center}

\begin{center}
{\large Raziel Carvajal-G{\'o}mez}\\
Universit{\'e} Catholique de Louvain\\
Louvain-la-Neuve, Belgium\\ 
\end{center}

\date{}

{\bf Keywords:} \emph{broadcast protocols, wireless networks, emergent overlays.}\\

%\section{Introduction}
Mobile ad-hoc networks (MANETs) do not rely on fixed infrastructures for their deployment.
%
Nodes in such a network have wireless transceivers with a limited transmission range and forward packets to communicate with each other in a peer-to-peer fashion.
%
Application scenarios such as vehicular networks, monitoring of zones in danger or the deployment of communication infrastructures in rural zones are suitable scenarios for MANETs where broadcast is a fundamental operation.
%
%\rcg{give concrete examples of systems where broadcast is important}
%

Broadcast protocols for MANETs tackle the problem of selecting a set of forwarding nodes that retransmit received packets to reach every node in the network.
%
Choosing a forwarding set is not an easy task, a small set reduces the overall energy consumption but may result in unreachable packets for other nodes.
%
On the contrarily, a bigger forwarding set increase the collisions and contentions in the network.
%
Mobility is another factor, if a forwarding set is not updated according to the velocity of nodes, obsolete choices of this set may cause lost of packets.
%
Additionally, mobility patterns may result in partitions and heterogeneous networks where the density of nodes varies within the area of communication.

%%
We are interested in designing adaptive approaches for broadcast protocols in MANETs.
%
In our context, adaptation refers to the ability of dynamically selecting and configuring the protocol use to broadcast.
%
Notice that the deployment of a broadcast protocol may take place either in the entire system, or more interestingly, just in a part of it.
%
To deal with the heterogeneous conditions of mobility and density, we are working in the design of adaptive approaches that converge to hybrid systems: overlays emerge in highly dense regions where nodes remains mostly static and, controlled flooding protocols are deployed in sparse regions where nodes move more rapidly.

%

In our vision, nodes discover (only with local information) the need of adaptation and perform a collective decision towards a switch to a different protocol.
%
Preliminary results of our simulations show that the energy cost tends to reduce with our adaptation approach, keeping a high coverage of nodes without overkilling the network bandwidth.

\newpage

\hypertarget{lastpage}{}
\end{document}

